\subsection{Partial and Total Derivatives}

% **Role:** calculus as local linearization plus global geometric integration.
%
% Linear approximation, Jacobians, partial/directional derivatives, chain rule.

Several-variable differentiation studies best linear approximations and local geometry, using linear maps, operator norms, determinants, invertibility, and metric estimates.

\begin{definition} \label{def:several-variable-differentiation-core-source-todo}\label{def:total-derivative-rn-rm} \hlt{Total Derivative}\\
Let $E\subseteq\R^n$ be open, $\f:E\to\R^m$, $\x\in E$, and $\LL(\R^n,\R^m)$ denote space of linear maps (Definition \ref{def:linear-map}).\\
If $\exists A\in\LL(\R^n,\R^m)$ s.t. $\lim\limits_{\h\to\ze}\frac{\norm{\f(\x+\h)-\f(\x)-A\h}_2}{\norm{\h}_2}=0$, then $\f$ is \hlt{differentiable at $\x$}, and $A$ denoted $D\f(\x)$ or $\f'(\x)$.\\
Equivalently, $\f(\x+\h)-\f(\x)=D\f(\x)\h+\mathbf{r}(\h)$, and $\frac{\norm{\mathbf{r}(\h)}_2}{\norm{\h}_2}\to0$.\\
If this holds for every $\x\in E$, then $\f$ is differentiable on $E$.
\end{definition}

\begin{perspective} \label{per:total-derivative-linear-map-first} \hlt{Derivative First, Matrix Second}\\
The total derivative is a linear map first; the Jacobian is its matrix in chosen bases. This keeps coordinate formulas subordinate to the intrinsic estimate $\norm{A\h}_2\leq\normop{A}\norm{\h}_2$ from Definition \ref{def:operator-norm-rn-rm}.
\end{perspective}

\begin{figure}[H]
  \centering
  \begin{tikzpicture}[
  scale=0.88,
  >=Stealth,
  every node/.style={font=\scriptsize},
  axis/.style={->, gray!70, line width=0.45pt},
  vector/.style={->, black, line width=0.75pt},
  linear/.style={->, green!45!black, line width=0.85pt},
  error/.style={->, black!65, dashed, dash pattern=on 3pt off 2pt, line width=0.55pt},
  point/.style={circle, fill=black, inner sep=1.15pt},
  approxpoint/.style={circle, fill=green!45!black, inner sep=1.15pt}
]
  \begin{scope}
    \draw[axis] (-0.25,0) -- (3.0,0) node[right] {$x$};
    \draw[axis] (0,-0.2) -- (0,2.1) node[above] {$y$};

    \coordinate (X) at (0.85,0.65);
    \coordinate (XH) at (2.25,1.25);

    \node[point] at (X) {};
    \node[point] at (XH) {};
    \draw[vector] (X) -- node[above left=-1pt] {$\h$} (XH);

    \node[below left=1pt] at (X) {$\x$};
    \node[right] at (XH) {$\x+\h$};
    \node[above] at (1.35,2.25) {domain};
  \end{scope}

  \begin{scope}[xshift=4.2cm]
    \draw[axis] (-0.25,0) -- (3.45,0) node[right] {$x$};
    \draw[axis] (0,-0.2) -- (0,2.35) node[above] {$y$};

    \coordinate (FX) at (0.65,0.72);
    \coordinate (LIN) at (2.18,1.45);
    \coordinate (ACT) at (2.98,1.95);

    \node[point] at (FX) {};
    \node[approxpoint] at (LIN) {};
    \node[point] at (ACT) {};

    \draw[linear] (FX) -- node[below right=-1pt] {$D\f(\x)\h$} (LIN);
    \draw[error] (LIN) -- node[below right=-1pt] {$\mathbf r(\h)$} (ACT);

    \node[left] at (FX) {$\f(\x)$};
    \node[above left=2pt] at (LIN) {$\f(\x)+D\f(\x)\h$};
    \node[above right=1pt] at (ACT) {$\f(\x+\h)$};
    \node[above] at (1.6,2.25) {codomain};
  \end{scope}
\end{tikzpicture}

  \caption{$D\f(\x)\h$ is the linear part; $\mathbf{r}(\h)=o(\norm{\h}_2)$ is the error.}
  \label{fig:total-derivative-linear-approximation}
\end{figure}

\begin{theorem} \label{thm:total-derivative-unique-and-continuous}\label{thm:several-variable-differentiation-main-source-todo} \hlt{Basic Consequences of Total Differentiability}\\
If $\f:E\to\R^m$ is differentiable at $\x\in E$, then $D\f(\x)$ is unique and $\f$ is continuous at $\x$.
\begin{proof}
If $A_1,A_2$ both satisfy Definition \ref{def:total-derivative-rn-rm}, set $B=A_1-A_2$. Then $\norm{B\h}_2/\norm{\h}_2\to0$ as $\h\to\ze$.\\
For fixed $\h\neq\ze$, applying this to $t\h$ and using linearity makes the quotient independent of $t$, so $B\h=\ze$.\\
Hence $B=0$. For continuity, write $\f(\x+\h)-\f(\x)=D\f(\x)\h+\mathbf{r}(\h)$.\\
The operator-norm estimate gives $\norm{D\f(\x)\h}_2\leq\normop{D\f(\x)}\norm{\h}_2$, and remainder is $o(\norm{\h}_2)$, so $\f(\x+\h)\to\f(\x)$.\\
\hlt{Structural Commentary}: The derivative is a rigid first-order linear part.\\
\hlt{Key Insight}: A remainder smaller than $\norm{\h}$ cannot hide a second linear map.
\end{proof}
\end{theorem}

\begin{example} \label{ex:derivative-of-linear-map} \hlt{Derivative of a Linear Map}\\
If $A\in\LL(\R^n,\R^m)$ and $\f(\x)=A\x$, then $D\f(\x)=A$ for every $\x\in\R^n$.\\
The remainder in Definition \ref{def:total-derivative-rn-rm} is identically zero.
\end{example}

\begin{theorem} \label{thm:chain-rule-several-variables} \hlt{Chain Rule}\\
If $\f:E\subseteq\R^n\to\R^m$ is differentiable at $\x_0$, $\g:G\subseteq\R^m\to\R^k$ is differentiable at $\f(\x_0)$, and $\f(E)\subseteq G$, then $\g\circ\f$ is differentiable at $\x_0$ and $D(\g\circ\f)(\x_0)=D\g(\f(\x_0))\,D\f(\x_0)$.
\begin{proof}
Let $A=D\f(\x_0)$, $B=D\g(\f(\x_0))$, $\mathbf{u}(\h)=\f(\x_0+\h)-\f(\x_0)-A\h$, and $\mathbf{v}(\mathbf{k})=\g(\f(\x_0)+\mathbf{k})-\g(\f(\x_0))-B\mathbf{k}$.\\
For $\mathbf{k}=\f(\x_0+\h)-\f(\x_0)$, differentiability of $\f$ gives $\mathbf{k}\to\ze$ and $(\g\circ\f)(\x_0+\h)-(\g\circ\f)(\x_0)-BA\h =B\mathbf{u}(\h)+\mathbf{v}(\mathbf{k})$, where $\norm{\mathbf{u}(\h)}_2=o(\norm{\h}_2)$ and $\norm{\mathbf{v}(\mathbf{k})}_2=o(\norm{\mathbf{k}}_2)=o(\norm{\h}_2)$. Divide by $\norm{\h}_2$.\\
\hlt{Structural Commentary}: The first-order parts compose as linear maps.\\
\hlt{Key Insight}: The output increment of $\f$ is small enough to feed into the remainder estimate for $\g$.
\end{proof}
\end{theorem}

\begin{theorem} \label{thm:several-variable-componentwise-differentiability} \hlt{Componentwise Differentiability}\\
Let $\f=(f_1,\ldots,f_m):E\subseteq\R^n\to\R^m$ and let $\pi_i:\R^m\to\R$ be the $i$th coordinate projection. Then $\f$ is differentiable at $\x\in E$ iff each $f_i$ is differentiable at $\x$, and $D f_i(\x)=\pi_i\circ D\f(\x)$ when these derivatives exist.
\begin{proof}
\begin{enumerate}[label=\roman*.]
\setlength{\itemsep}{0pt}
\item[$(\Rightarrow)$] If $\f$ is differentiable, write $\f(\x+\h)-\f(\x)=D\f(\x)\h+\mathbf r(\h)$ with $\norm{\mathbf r(\h)}_2=o(\norm{\h}_2)$.\\
Applying $\pi_i$ gives $f_i(\x+\h)-f_i(\x)=\pi_i(D\f(\x)\h)+\pi_i(\mathbf r(\h))$ and $\abs{\pi_i(\mathbf r(\h))}\leq\norm{\mathbf r(\h)}_2=o(\norm{\h}_2)$.\\
Thus $D f_i(\x)=\pi_i\circ D\f(\x)$.
\item[$(\Leftarrow)$] Suppose $D f_i(\x)=A_i\in\LL(\R^n,\R)$ for each $i$. Define $A\in\LL(\R^n,\R^m)$ by $A\h=(A_1\h,\ldots,A_m\h)$.\\
With $r_i(\h)=f_i(\x+\h)-f_i(\x)-A_i\h$, each $r_i(\h)=o(\norm{\h}_2)$.\\
So $\mathbf r(\h)=(r_1(\h),\ldots,r_m(\h))$ satisfies $\norm{\mathbf r(\h)}_2\leq\sum_i\abs{r_i(\h)}=o(\norm{\h}_2)$. Definition \ref{def:total-derivative-rn-rm} gives $D\f(\x)=A$.
\end{enumerate}
\hlt{Structural Commentary}: Component derivatives assemble into one linear map.\\
\hlt{Key Insight}: Components test the derivative; the total derivative packages the tests together.
\end{proof}
\end{theorem}

\begin{definition} \label{def:partials-gradient-directional-derivatives} \hlt{Partial, Gradient, and Directional Derivatives}\\
Let $E\subseteq\R^n$ be open, $\f=(f_1,\ldots,f_m):E\to\R^m$, and let $\mathbf{e}_1,\ldots,\mathbf{e}_n$ be the standard basis of $\R^n$.\\
If indicated limits exist, \hlt{partial derivative} is $(D_jf_i)(\x)\defeq\lim_{t\to0}\frac{f_i(\x+t\mathbf{e}_j)-f_i(\x)}{t}$ for $1\leq i\leq m,\ 1\leq j\leq n$.\\
\hlt{Directional derivative} is $(D_{\mathbf{u}}f)(\x)\defeq\lim_{t\to0}\frac{f(\x+t\mathbf{u})-f(\x)}{t}$ for $\mathbf{u}\in\R^n$; when only the direction is meant, one often assumes $\norm{\mathbf{u}}_2=1$. If $m=1$ and all partials exist, the \hlt{gradient} is $(\nabla f)(\x)\defeq \sum\limits_{j=1}^{n}(D_jf)(\x)\mathbf{e}_j$.
\end{definition}

\begin{perspective} \label{per:partials-are-coordinate-tests} \hlt{Partials Are Coordinate Tests}\\
Partials and directional derivatives test selected lines; the total derivative controls all small increments by one linear map. Theorem \ref{thm:total-derivative-matrix-partials} extracts partials, and Theorem \ref{thm:c1-iff-continuous-partials} gives a continuity-based converse.
\end{perspective}

\begin{theorem} \label{thm:total-derivative-matrix-partials} \hlt{Matrix of the Total Derivative}\\
If $\f:E\to\R^m$ is differentiable at $\x$, then every $(D_jf_i)(\x)$ exists and, in standard bases,
\begin{align}
[D\f(\x)] =
\begin{bmatrix}
(D_1f_1)(\x) & \cdots & (D_nf_1)(\x)\\
\vdots & & \vdots\\
(D_1f_m)(\x) & \cdots & (D_nf_m)(\x)
\end{bmatrix} \nonumber
\end{align}
If $m=1$ and $f$ is differentiable at $\x$, then $(D_{\mathbf{u}}f)(\x)=(\nabla f)(\x)\cdot\mathbf{u}$ for every $\mathbf{u}\in\R^n$.
\begin{proof}
Set $\h=t\mathbf{e}_j$ in Definition \ref{def:total-derivative-rn-rm}; limit is $D\f(\x)\mathbf{e}_j$, so components form column $j$ (Proposition \ref{prop:matrix-representation-action}).\\
For $m=1$, set $\h=t\mathbf{u}$, identify linear functional with dot product against $\nabla f(\x)$ (Theorem \ref{thm:riesz-representation-finite-dimensional}).\\
\hlt{Structural Commentary}: Partials are coordinate columns of $D\f(\x)$.\\
\hlt{Key Insight}: Differentiability controls all directions at once; partial derivatives record its coordinate columns.
\end{proof}
\end{theorem}

\begin{remark} \label{rem:several-variable-row-column-conventions} \hlt{Row and Column Conventions}\\
With the coordinate-column convention from Proposition \ref{prop:matrix-representation-action}, the Jacobian matrix above is the standard-basis matrix of the linear map $D\f(\x)$, and its $j$th column is $D\f(\x)\mathbf{e}_j$.
\end{remark}

